% ================== Chapitre 1 : objectifs, perimetre et exigences ==========
\chapter{Objectifs, périmètre et exigences couvertes}
\label{ch:objectifs}

\section{Objectif de la tâche et définition du service attendu}

La tâche T005 consiste à mettre en place le service DNS de GANDAL. Dans le
découpage du suivi d'équipe, elle se situe après le gel du plan d'adressage et
du modèle IPAM, et en parallèle du DHCP haute disponibilité et de la
synchronisation temporelle. Le service attendu n'est pas un simple serveur de
noms : le document 03 lui assigne quatre fonctions distinctes.

\begin{itemize}
  \item \textbf{Publier l'espace de nommage interne du datacenter}, c'est-à-dire
    les noms d'infrastructure et les noms des machines de chaque tenant, avec
    les résolutions inverses correspondantes.
  \item \textbf{Fournir une résolution récursive contrôlée} aux machines
    autorisées, pour les noms internes comme pour les noms publics.
  \item \textbf{Garantir la confidentialité des noms entre domaines} : un tenant
    ne doit lire ni la zone d'infrastructure ni la zone d'un autre tenant.
  \item \textbf{S'insérer dans le cycle de vie d'une allocation} : la publication
    d'un enregistrement est une étape de la machine d'états de création, entre
    l'attribution du bail et l'exposition de l'état \opt{READY}.
\end{itemize}

Mettre en place le DNS signifie donc arrêter une architecture, un espace de
nommage, des règles d'accès et un mécanisme de mise à jour, puis produire les
configurations et procédures qui permettent de les appliquer de manière
reproductible. C'est l'objet des chapitres \ref{ch:conception} à
\ref{ch:cycle}.

\begin{encadre}
\titreencadre{Règle de restitution du corpus}
Tous les documents de cadrage de GANDAL précisent que la règle est d'effectuer
une simulation au préalable et qu'aucun déploiement ni essai n'est présenté
comme effectué. Ce dossier respecte cette règle : il fournit des configurations
et des procédures complètes, mais les contrôles du chapitre~\ref{ch:recette}
portent le statut \statut{non exécuté}. Produire les preuves correspondantes
relève de la tâche T009, qui n'est pas traitée ici.
\end{encadre}

\section{Contexte repris du corpus et données d'entrée}

Les choix qui suivent sont contraints par des décisions déjà arrêtées. Elles ne
sont pas rediscutées ici ; elles sont rappelées parce qu'elles déterminent
directement la conception du service.

\begin{longtable}{@{}R{0.31}R{0.52}R{0.17}@{}}
\caption{Données d'entrée de la tâche T005 et leurs conséquences directes}
\label{tab:entrees}\\
\hautfilet
\textbf{Contrainte ou décision antérieure} &
\textbf{Conséquence sur le service DNS} & \textbf{Source}\\
\medfilet
\endfirsthead
\hautfilet
\textbf{Contrainte ou décision antérieure} &
\textbf{Conséquence sur le service DNS} & \textbf{Source}\\
\medfilet
\endhead
\basfilet
\endfoot
Trois hôtes XCP-ng, une NIC par hôte, un commutateur L2 et un routeur non
  manageables &
  La redondance logicielle ne peut couvrir que la perte d'une machine virtuelle
  ou d'un hôte. SW1 et R1 restent des dépendances communes assumées. &
  Document 01, 2.1\\
Bloc \ip{10.30.0.0/16}, VNet Services en \ip{10.30.4.0/24}, DNS autoritatif en
  \ip{.10} et \ip{.11}, DNS récursif en \ip{.12} et \ip{.13} &
  Les adresses des quatre instances sont imposées. Ce dossier les consomme sans
  redéfinition locale. & Document 02, tableau 2.5\\
Tenant $n$ associé au préfixe \ip{10.30.(63+n).0/24}, pour
  $1\leqslant n\leqslant 50$ &
  La correspondance entre un tenant, sa zone directe et sa zone inverse se
  déduit arithmétiquement, ce qui rend la génération des zones et la politique
  de vue déterministes. & Document 02, 2.4\\
VXLAN sur IPv4 avec MTU externe de 1\,500 octets &
  Au plus 1\,450 octets utiles pour le paquet IP interne. Le dimensionnement
  des réponses DNS en doit tenir compte. & Document 02, 2.5\\
PowerDNS Authoritative complété de PowerDNS Recursor &
  Choix proposé par la comparaison du document 03 ; BIND reste l'alternative
  documentée du corpus. & Document 03, tableau 1.1\\
Aucun VNet surveillé n'initie de connexion vers Monitoring &
  Les journaux et métriques du DNS sont conservés localement et lus en pull.
  Aucun envoi direct n'est configuré. & Document 03, 2.7\\
Deux instances de chaque service, sur des hôtes distincts &
  Objectif de continuité logique, sans prétendre à l'élimination des points
  uniques de défaillance. & \exig{net-nfr-001}\\
\end{longtable}

\section{Registre des exigences couvertes}

Le tableau \ref{tab:tracabilite} relie chaque exigence du registre des services
réseau au mécanisme retenu dans ce dossier et au contrôle de recette qui
permettra de la vérifier. Les identifiants de contrôle \opt{C01} à \opt{C15}
sont définis au chapitre~\ref{ch:recette}.

\begin{longtable}{@{}R{0.16}R{0.18}R{0.50}R{0.16}@{}}
\caption{Traçabilité entre les exigences du registre et les mécanismes de ce
  dossier}\label{tab:tracabilite}\\
\hautfilet
\textbf{Exigence} & \textbf{Objet} & \textbf{Mécanisme retenu et section} &
\textbf{Contrôle}\\
\medfilet
\endfirsthead
\hautfilet
\textbf{Exigence} & \textbf{Objet} & \textbf{Mécanisme retenu et section} &
\textbf{Contrôle}\\
\medfilet
\endhead
\basfilet
\endfoot
NET-FR-004 & Paramètres DNS par tenant & Objet profil DNS porté par l'IPAM et
  consommé par DHCP et Templates (\ref{sec:profil}) & C03\\
NET-FR-006 & DNS interne et mises à jour & Zones internes autoritatives et
  publication par l'agent via l'API (\ref{sec:nommage}, \ref{sec:publication})
  & C01, C10\\
NET-FR-007 & Récursion DNS autorisée & \opt{allow-from} restreint aux préfixes
  du plan, sortie 53 explicite (\ref{sec:recursif}) & C07\\
NET-FR-008 & Confidentialité des noms entre domaines & Politique de vue par
  adresse source, transferts fermés, autoritatif non exposé
  (\ref{sec:vue}, \ref{sec:flux}) & C04, C05, C06\\
NET-FR-012 & Allocation et libération idempotentes & Modification par RRset avec
  lecture préalable et clé d'idempotence (\ref{sec:idempotence}) & C10\\
NET-FR-013 & Paramètres VM via Templates & Profil DNS transmis par le contrat
  IF-NS-06 (\ref{sec:profil}) & C03\\
NET-FR-014 & Métriques services & Point de métriques en lecture seule,
  interrogé en pull (\ref{sec:journaux}) & C14\\
NET-FR-016 & API administratives chiffrées & API autoritative liée à la boucle
  locale derrière une terminaison TLS (\ref{sec:api}) & C10\\
NET-FR-017 & DNS accessible selon profil VPN & Vue administration pour le pool
  VPN admin, refus par défaut pour le pool utilisateur (\ref{sec:vpn}) & C06\\
NET-FR-018 & Journaux horodatés & Journalisation dnstap horodatée en temps
  universel, conservée localement (\ref{sec:journaux}) & C14\\
NET-FR-019 & Gestion de l'infrastructure isolée & Les VNets tenant ne joignent
  ni le plan de gestion ni les autoritatifs (\ref{sec:flux}) & C08\\
NET-FR-020 & Simulation avant déploiement & Statut \statut{non exécuté}
  maintenu ; maquette préalable exigée (chapitre~\ref{ch:limites}) & aucun\\
NET-NFR-001 & Deux instances DNS & Deux autoritatifs et deux récursifs répartis
  sur les trois hôtes (\ref{sec:placement}) & C09\\
NET-NFR-003 & Reconstruction en moins de 60 min & Export des zones et procédure
  de restauration (\ref{sec:sauvegarde}) ; mesure confiée à T009 & aucun\\
NET-NFR-005 & Protection des communications administratives & TLS sur l'API,
  TSIG sur les transferts (\ref{sec:zones}, \ref{sec:api}) & C04\\
NET-NFR-006 & Rejeu déterministe et idempotent & Opérations convergentes sur
  l'état désiré des RRsets (\ref{sec:idempotence}) & C10\\
NET-NFR-007 & Secrets absents des dépôts en clair & Clés API et TSIG injectées
  au démarrage ; empreinte en configuration (\ref{sec:sauvegarde}) & aucun\\
IF-NS-03 & DNS vers VNet et SDN : profils et ACL par tenant & Étiquette de vue
  dérivée du préfixe tenant (\ref{sec:vue}) & C05\\
IF-NS-06 & Agent vers Templates : réseau, DNS, MTU & Contenu normalisé du
  profil DNS (\ref{sec:profil}) & C03\\
IF-NS-08 & Network vers Monitoring : lecture en pull & Aucune initiation
  sortante depuis les VM DNS (\ref{sec:journaux}) & C14\\
\end{longtable}

\section{Limites de périmètre et articulation avec les autres tâches}

Trois tâches voisines touchent au même cycle de vie sans relever de T005. Les
points de contact sont explicités pour éviter une double définition.

\begin{longtable}{@{}R{0.30}R{0.70}@{}}
\caption{Frontières de la tâche T005}\label{tab:frontieres}\\
\hautfilet
\textbf{Tâche voisine} & \textbf{Point de contact et répartition}\\
\medfilet
\endfirsthead
\hautfilet
\textbf{Tâche voisine} & \textbf{Point de contact et répartition}\\
\medfilet
\endhead
\basfilet
\endfoot
T004, DHCP haute disponibilité \newline (B.~Heudep Djandja) &
  T005 fournit les adresses des résolveurs et le nom de domaine à distribuer
  dans les options 6, 15 et 119. T005 ne définit ni les scopes, ni le mode de
  haute disponibilité de Kea, ni la durée de bail. La valeur de bail est ici
  une hypothèse de travail, signalée comme telle.\\
T006, synchronisation temporelle \newline (B.~Heudep Djandja) &
  La validation DNSSEC et la journalisation corrélable supposent une horloge
  correcte. T005 énonce cette dépendance et la contrainte d'amorçage à froid,
  mais ne configure pas chrony.\\
T007, agent réseau et API de provisionnement \newline (P.~Pangui Pianne) &
  T005 définit le contrat que l'agent doit respecter pour publier un
  enregistrement et le comportement attendu en cas d'échec. L'implémentation de
  l'agent ne relève pas de ce dossier.\\
T009, tests, validation, reprise et documentation &
  Seconde tâche attribuée au même responsable. Elle dépend des résultats des
  lots précédents et n'est pas traitée ici. Le chapitre~\ref{ch:recette} lui
  transmet des critères d'acceptation déjà formulés, qui devront être exécutés
  et archivés dans ce cadre.\\
\end{longtable}
